\newif\ifglobal

%%%%%%%%%%%%%%%%%%%%%%%
%%% Compile to view %%%
%%%%%%%%%%%%%%%%%%%%%%%

%%%%%%%%%%%%%%%%%%%%%%%%%%%%%%%%%%%%%%%%%%%%%%%%%%%
%%% Readme:                                     %%%
%%% https://www.overleaf.com/read/bwdjvfjksvjy/ %%%
%%%%%%%%%%%%%%%%%%%%%%%%%%%%%%%%%%%%%%%%%%%%%%%%%%%

% >>> M?: marks a module setting number or important notice

% >>> M4: Set {./relative-path} to external files for this module <<<
\def \modulefiles {./08-RESCLK}
% To add an external file, use:
% (Replace '---' with actual filename)
% '\input{\modulefiles/tables/---.tex}' for tables
% '\includegraphics{\modulefiles/figures/---}' for figures
% '\subfile{\modulefiles/---}' for register files


% >>> M1: This module is in global project? <<<
% 'YES' - uncomment; 'NO' - comment out
\globaltrue

\ifglobal
    \documentclass[main\_\_EN.tex]{subfiles}

\else
    \documentclass[openany, 10pt]{book}

    % >>> M2: In ./00-shared/preamble-trm-module, also find
    %         settings for visibility of todo notes, watermark <<<
    \usepackage{preamble-trm-module}

    % >>> M3: Temporary labels in English,
    %         you can add your own labels there <<<
    \myexternaldocument{./00-shared/temp-labels-en}

\fi %\ifglobal


\setcounter{secnumdepth}{4}


% >>> M5: If this module needs any updates in
% './00-shared/preamble-trm-module.sty'
% (include additional package, etc.), contact your module PM
% or create the file './00-shared/changelog.tex' make the
% necessary updates yourself and document them in this file <<<

\raggedright


\begin{document}


% Sets language into which pre-defined bits of text are translated
\selectlanguage{english}

% Do not delete this block
\ifglobal
% Add the following front matter if
% this module is outside of global project
\else
    \InputIfFileExists{readme}{}{}
    {\let\clearpage\relax \listoftodos}
    {\let\clearpage\relax \listofdonetodos}
    \tableofcontents
    %\listoftables
    %\listoffigures
\fi


%%%%%%%%%%%%%%%%%%%%%%%%%
%%% TRM Module Begins %%%
%%%%%%%%%%%%%%%%%%%%%%%%%

\hypertarget{resclk}{}
\chapter{Reset and Clock}
\label{mod:resclk}


\section{Reset}
\subsection{Overview}
\chipname{} provides four types of reset that occur at different levels, namely CPU Reset, Core Reset, System Reset, and Chip Reset. All reset types mentioned above (except Chip Reset) preserve the data stored in internal memory. Figure  \ref{figure:Reset} shows the scopes of affected subsystems by each type of reset.

\subsection{Architectural Overview}
\begin{figure}[H]
    \centering
    \includegraphics[width=0.8\textwidth]{\modulefiles/figures/reset-overview.pdf}
    \caption{Reset Types}
    \label{figure:Reset}
\end{figure}
\chipname{}'s Digital System can be divided into two parts: \hyperref[mod:riscvcpu]{High Performance System (HP system)} that includes Digital Core, Wireless Circuit, and HP SRAM, and \hyperref[mod:lowpowm]{Low Power System (LP system)} that only includes some low-power peripherals and LP SRAM. See Figure \ref{figure:Reset} for details (note that HP SRAM and LP SRAM would not be reset, so they are not shown in the figure).

\subsection{Features}
\label{sec:resclk-features}
\begin{itemize}
    \item Four reset types:
    \begin{itemize}
    \item CPU Reset: resets CPU core. Once such a reset is released, the instructions from the CPU reset vector (0x40000000)  will be executed.
    \item Core Reset: resets the whole digital system except LP system, including CPU, peripherals, digital GPIOs, Bluetooth$^®$ LE, and 802.15.4.
    \item System Reset: resets the whole digital system, including LP system.
    \item Chip Reset: resets the whole chip.
    \end{itemize}
    \item Software reset and hardware reset:
    \begin{itemize}
        \item Software Reset: triggered via software by configuring the corresponding registers of CPU, see Chapter \ref{mod:lowpowm} \textit{\nameref{mod:lowpowm}}.
        \item Hardware Reset: triggered directly by the hardware.
    \end{itemize}
\end{itemize}

\subsection{Functional Description}
CPU will be reset immediately when any type of reset above occurs. Users can retrieve reset source codes by reading  \hyperref[fielddesc:RTCCLKRSTRESETCAUSE]{RTC\_CLKRST\_RESET\_CAUSE} after the reset is released.
Table \ref{tab:resclk-reset-src} lists possible reset sources and the types of reset they trigger.


\begin{table}[H]
    \centering
    \caption{Reset Source}
    \label{tab:resclk-reset-src}
    \begin{threeparttable}
    \begin{tabular}{ |l|m{3.5cm}|m{1.8cm}|m{9.2cm}| }
    \hline
\rowcolor{lightgray}
\textbf{Code} & \textbf{Source} & \textbf{Reset Type} &   \textbf{Note} \\\hline
        0x01 & Chip reset\tnote{1} & Chip Reset &---\\
        \hline
        0x0F & Brown-out system reset & Chip Reset or System Reset & Triggered by brown-out detector\tnote{2}\\
        \hline
        0x10 & RWDT system reset & System Reset & See Chapter \ref{mod:wdt} \textit{\nameref{mod:wdt}}
        \\\hline
        0x12 & Super Watchdog reset & System Reset & See Chapter \ref{mod:wdt} \textit{\nameref{mod:wdt}}
        \\\hline
        0x03 & Software system reset & Core Reset & Triggered by configuring \hyperref[fielddesc:LPAONHPSYSSWRESET]{LP\_AON\_HPSYS\_SW\_RESET}\\
        \hline
        0x05 & Deep-sleep reset & Core Reset & See Chapter \ref{mod:lowpowm} \textit{\nameref{mod:lowpowm}}
        \\\hline
        0x07 & MWDT0 core reset & Core Reset & See Chapter \ref{mod:wdt} \textit{\nameref{mod:wdt}}
        \\\hline
        0x08 & MWDT1 core reset & Core Reset & See Chapter \ref{mod:wdt} \textit{\nameref{mod:wdt}}
        \\\hline
        0x09 & RWDT core reset & Core Reset & See Chapter \ref{mod:wdt} \textit{\nameref{mod:wdt}}
        \\\hline
        0x14 & eFuse reset & Core Reset & Triggered by eFuse CRC error
        \\\hline
        0x15 & USB (UART) reset & Core Reset & Triggered when external USB host sends a specific command to the serial interface of USB Serial/JTAG Controller. See Chapter \ref{mod:usbserialjtag} \textit{\nameref{mod:usbserialjtag}} \\\hline
        0x16 & USB (JTAG) reset & Core Reset & Triggered when external USB host sends a specific command to the JTAG interface of USB Serial/JTAG Controller. See Chapter \ref{mod:usbserialjtag} \textit{\nameref{mod:usbserialjtag}} \\\hline
        0x0B & MWDT0 CPU reset & CPU Reset & See Chapter \ref{mod:wdt} \textit{\nameref{mod:wdt}}
        \\\hline
         0x0C & Software CPU reset & CPU Reset & Triggered by configuring \hyperref[fielddesc:LPAONCPUCORE0SWRESET]{LP\_AON\_CPU\_CORE0\_SW\_RESET}\\
        \hline
        0x0D & RWDT CPU reset & CPU Reset & See Chapter \ref{mod:wdt} \textit{\nameref{mod:wdt}}
        \\\hline
        0x11 & MWDT1 CPU reset & CPU Reset & See Chapter \ref{mod:wdt} \textit{\nameref{mod:wdt}}
        \\\hline
        0x17 & Voltage glitch reset & System Reset & Triggered by power supply glitch attack
        \\\hline
        0x18 & JTAG CPU reset & CPU Reset & Triggered when a "JDB Resetting CPU" instruction is received
        \\\hline
\end{tabular}
        \begin{tablenotes}
            \item[1] Chip Reset can be triggered by the following sources:
                \begin{itemize}
                \item Triggered by chip power-up.
                \item Triggered by brown-out detector.
                \end{itemize}
            \item[2] Once brown-out status is detected, the detector will trigger System Reset or Chip Reset, depending on the register configuration. See Chapter  \ref{mod:lowpowm} \textit{\nameref{mod:lowpowm}}.
        \end{tablenotes}
     \end{threeparttable}
    \end{table}


\subsection{Peripheral Reset}

Peripherals can be reset individually by configuring corresponding registers, or globally by core reset, system reset, or chip reset. After releasing peripherals from reset, the reset release status registers can be read to determine the reset status. These registers turning to 1 indicates that peripherals have been released from reset and can operate properly.

The reset registers of \chipname{} peripherals are controlled by the Power/Clock/Reset (PCR) module. For reset registers of peripherals, see Section \ref{sec:resclk-reg-summ} \textit{\nameref{sec:resclk-reg-summ}}.
\section{Clock}

\subsection{Overview}

\chipname{} clocks are mainly sourced from oscillator (OSC), RC, and PLL circuits, and then processed by the dividers or selectors, which allows most functional modules to select their working clock according to their power consumption and performance requirements. Figure  \ref{figure:Clock} shows the system clock structure.

\subsection{Architectural Overview}

\begin{figure}[H]
    \centering
    \includegraphics[width=1.0\textwidth]{\modulefiles/figures/sys_clk_struct.pdf}
    \caption{System Clock}
    \label{figure:Clock}
\end{figure}


\subsection{Features}

\chipname{} clocks can be classified into two types depending on their frequencies:
\begin{itemize}
    \item High-performance (HP) clocks for devices working at a higher frequency, such as CPU and digital peripherals
    \begin{itemize}
        \item PLL\_F96M\_CLK (96 MHz): internal PLL clock. Its reference clock is XTAL\_CLK.
        \item PLL\_F64M\_CLK (64 MHz): internal PLL clock. Its reference clock is XTAL\_CLK.
        \item XTAL\_CLK (32 MHz): external crystal clock
    \end{itemize}
    \item Low-power (LP) clocks for low-power system and some peripherals working in low-power mode
    \begin{itemize}
        \item XTAL32K\_CLK (32 kHz): external crystal clock
	    \item RC\_FAST\_CLK (8 MHz by default): internal fast RC oscillator with adjustable frequency
	    \item RC\_SLOW\_CLK (130 kHz): internal slow RC oscillator
	    \item OSC\_SLOW\_CLK (32 kHz by default): external slow clock input through \hyperref[tab:iomuxgpio-iomax-pins-analog-functions]{XTAL\_32K\_P}. After configuring this GPIO, also configure the Hold function (see Chapter \ref{mod:iomuxgpio} \textit{\nameref{mod:iomuxgpio}} >  \ref{PadHold} \textit{\nameref{PadHold}})
        \item PLL\_LP\_CLK (8 MHz):  internal PLL clock, with reference clock XTAL32K\_CLK
    \end{itemize}
\end{itemize}

\subsection{Functional Description}

\subsubsection{HP System Clock}

As Figure \ref{figure:Clock} shows, CPU\_CLK is the master clock for CPU and its frequency can be as high as 96 MHz. Alternatively, CPU can run at lower frequencies, such as at 2 MHz, to achieve lower power consumption. CPU\_CLK shares the same clock sources with AHB\_CLK and APB\_CLK. Users can select from XTAL\_CLK, PLL\_96M\_CLK, PLL\_64M\_CLK, or RC\_FAST\_CLK as the clock source of CPU\_CLK by configuring  \hyperref[fielddesc:PCRSOCCLKSEL]{PCR\_SOC\_CLK\_SEL}. For details, see Table  \ref{tab:resclk-cpu-clk-src} and Table \ref{tab:resclk-cpu-clk-freq}. By default, the CPU clock is sourced from XTAL\_CLK with a divider of 1, i.e., the CPU clock frequency is 32 MHz.

\begin{longtable}[c]{|p{5cm}|p{5cm}| }
       \caption{CPU\_CLK Clock Source}
		\label{tab:resclk-cpu-clk-src} \\
	    \hline
        \rowcolor{lightgray}
        \textbf{\hyperref[fielddesc:PCRSOCCLKSEL]{PCR\_SOC\_CLK\_SEL}} & \textbf{CPU Clock Source} \\
        \hline
        0 & XTAL\_CLK \\
        \hline
        1 & PLL\_F96M\_CLK \\
        \hline
        2 & RC\_FAST\_CLK \\
        \hline
        3 & PLL\_F64M\_CLK \\
        \hline
\end{longtable}

\begin{table}[H]
    \centering
    \caption{Frequency of CPU\_CLK, AHB\_CLK and HP\_ROOT\_CLK}
    \label{tab:resclk-cpu-clk-freq}

    \begin{threeparttable}
    \begin{tabular}{ |p{2.3cm}|p{4.0cm}|p{8.5cm}|}
    \hline
        \rowcolor{lightgray}
        \textbf{Clock} & \textbf{Source} & \textbf{Frequency} \\\hline
		\multirow{4}{*} {HP\_ROOT\_CLK} & PLL\_F96M\_CLK & 96 MHz \\\cline{2-3}
		& PLL\_F64M\_CLK & 64 MHz \\\cline{2-3}
		& XTAL\_CLK & 32 MHz \\\cline{2-3}
		& RC\_FAST\_CLK & 8 MHz\\\hline
		{CPU\_CLK\tnote{1}} & HP\_ROOT\_CLK &  $f_{\textrm{HP\_ROOT\_CLK}}$/(\hyperref[fielddesc:PCRCPUDIVNUM]{PCR\_CPU\_DIV\_NUM}+1 )\\\hline
		{AHB\_CLK\tnote{2}} & HP\_ROOT\_CLK & $f_{\textrm{HP\_ROOT\_CLK}}$/(\hyperref[fielddesc:PCRAHBDIVNUM]{PCR\_AHB\_DIV\_NUM}+1 )\\\hline

    \end{tabular}

        \begin{tablenotes}
            \item[1] CPU\_CLK frequency must be larger than or equal to AHB frequency, and must be an integer multiple of AHB\_CLK frequency.
            \item[2] AHB\_CLK frequency can not exceed 32 MHz.
        \end{tablenotes}
     \end{threeparttable}
    \end{table}

\begin{tiplisting}
When selecting the clock source of HP\_ROOT\_CLK, or configuring the clock divisor for CPU\_CLK and AHB\_CLK, please also set \hyperref[fielddesc:PCRBUSCLOCKUPDATE]{PCR\_BUS\_CLOCK\_UPDATE} to apply the new configuration, and read \hyperref[fielddesc:PCRBUSCLOCKUPDATE]{PCR\_BUS\_CLOCK\_UPDATE} to see if new configuration takes effect.
\end{tiplisting}

As shown in \ref{figure:Clock}, to generate APB\_CLK, AHB\_CLK might be divided twice. The first division is compulsory. That is, AHB\_CLK is always divided by the divisor (\hyperref[fielddesc:PCRAPBDIVNUM]{PCR\_APB\_DIV\_NUM} + 1). The second division (also called automatic frequency reduction) is optional. When there is no request from the host in the chip to access peripheral registers, AHB\_CLK will be further divided by (\hyperref[fielddesc:PCRAPBDECREASEDIVNUM]{APB\_DECREASE\_DIV\_NUM} + 1) to lower power consumption. If the host initiates a request to access peripheral registers, APB\_CLK will be restored to the frequency after the first division.

Note that the chip's performance will degrade due to the automatic frequency reduction. This function can be disabled (already disabled by default) by configuring \hyperref[fielddesc:PCRAPBDECREASEDIVNUM]{APB\_DECREASE\_DIV\_NUM} to 0.

\subsubsection{LP System Clock}

The LP system can operate when most other clocks are disabled. LP system clocks include LP\_SLOW\_CLK and LP\_FAST\_CLK.

The clock sources for LP\_SLOW\_CLK and LP\_FAST\_CLK are low-frequency clocks:
\begin{itemize}
    \item LP\_SLOW\_CLK can be derived from:
    \begin{itemize}
        \item RC\_SLOW\_CLK
        \item XTAL32K\_CLK
        \item OSC\_SLOW\_CLK
    \end{itemize}
    \item LP\_FAST\_CLK can be derived from:
    \begin{itemize}
        \item 16 MHz XTAL\_D2\_CLK, which is XTAL\_CLK divided by 2
        \item RC\_FAST\_CLK
        \item PLL\_LP\_CLK
    \end{itemize}
\end{itemize}

The clock source of LP\_DYN\_SLOW\_CLK is LP\_SLOW\_CLK. There is no frequency change from the clock source.

The clock source of LP\_DYN\_FAST\_CLK depends on the chip's power mode (see Chapter \ref{mod:lowpowm} \textit{\nameref{mod:lowpowm}}).
\begin{itemize}
    \item Select LP\_FAST\_CLK as its clock source in Active and Modem-sleep mode
    \item Select LP\_SLOW\_CLK as its clock source in Light-sleep and Deep-sleep mode
\end{itemize}

\subsubsection{Peripheral Clocks}
Table \ref{tab:resclk-periph-clks-root-hp}, Table \ref{tab:resclk-periph-clks-hp}, Table \ref{tab:resclk-periph-clks-root-lp}, and Table \ref{tab:resclk-periph-clks-lp} list the derived HP/LP clocks sources and HP/LP clocks for each peripheral.

\begin{landscape}

\setlength\tabcolsep{1pt}

\begin{tiny}
\begin{longtable}[c]{|m{2cm}|c|c|c|c|c|c|c|c|c|C{1.5cm}|c|c|c|c|c|c|c|c|c| }

 \caption{Derived HP Clock Source}
    \label{tab:resclk-periph-clks-root-hp} \\
  \hline
    \rowcolor{lightgray}
    & \multicolumn{3}{c|}{\textbf{Source Clock}} & \textbf{Derived Clock} & \multicolumn{5}{c|}{\textbf{Source Clock}} & \multicolumn{9}{c|}{\textbf{Derived Clock}} & \textbf{Source Clock} \\ \cline{2-20}
    \rowcolor{lightgray}
    & & \multicolumn{3}{c|}{\textbf{PLL\_CLK}} &  &  &  &  &  & \textbf{HP\_ROOT\_CLK} &  &  &  &  &  &  & &  & \multirow{2}{*}{\textbf{}} \\ \cline{3-5}
    \rowcolor{lightgray}
    \multirow{-3}{*}{\textbf{Derived Clock}} & \multirow{-2}{*}{\textbf{\makecell{XTAL\_CLK\\32 MHz}}} & \textbf{\makecell{PLL\_F96M\_CLK\\96 MHz}} & \textbf{\makecell{PLL\_F64M\_CLK\\64 MHz}} & \textbf{\makecell{PLL\_F48M\_CLK\\48 MHz}} & \multirow{-2}{*}{\textbf{\makecell{RC\_FAST\_CLK\\8 MHz}}} & \multirow{-2}{*}{\textbf{\makecell{RC\_SLOW\_CLK\\130 kHz}}} & \multirow{-2}{*}{\textbf{\makecell{OSC\_SLOW\_CLK\\32 kHz}}} & \multirow{-2}{*}{\textbf{\makecell{XTAL32K\_CLK\\32 kHz}}} & \multirow{-2}{*}{\textbf{\makecell{PLL\_LP\_CLK\\8 MHz}}} & \textbf{96 MHz/64 MHz/32 MHz/8 MHz} & \multirow{-2}{*}{\textbf{CRYPTO\_CLK}} & \multirow{-2}{*}{\textbf{APB\_CLK}} &  \multirow{-2}{*}{\textbf{AHB\_CLK}} & \multirow{-2}{*}{\textbf{CPU\_CLK}} & \multirow{-2}{*}{\textbf{\makecell{LP\_DYN\_\\FAST\_CLK}}}  & \multirow{-2}{*}{\textbf{\makecell{LP\_DYN\_\\SLOW\_CLK}}} & \multirow{-2}{*}{\textbf{\makecell{XTAL\_D2\_CLK\\16 MHz}}} &  \multirow{-2}{*}{\textbf{LP\_FAST\_CLK}} & \multirow{-2}{*}{\textbf{\makecell{Clock\\from IO}}} \\ \hline

        \endfirsthead

    \multicolumn{20}{c}{\bfseries \tablename\ \thetable{} -- Continued from the previous page...}\\\hline
    \rowcolor{lightgray}
    & \multicolumn{3}{c|}{\textbf{Source Clock}} & \textbf{Derived Clock} & \multicolumn{5}{c|}{\textbf{Source Clock}} & \multicolumn{9}{c|}{\textbf{Derived Clock}} & \textbf{Source Clock} \\ \cline{2-20}
    \rowcolor{lightgray}
    & & \multicolumn{3}{c|}{\textbf{PLL\_CLK}} &  &  &  &  &  & \textbf{HP\_ROOT\_CLK} &  &  &  &  &  &  & &  & \multirow{2}{*}{\textbf{}} \\ \cline{3-5}
    \rowcolor{lightgray}
    \multirow{-3}{*}{\textbf{Derived Clock}} & \multirow{-2}{*}{\textbf{\makecell{XTAL\_CLK\\32 MHz}}} & \textbf{\makecell{PLL\_F96M\_CLK\\96 MHz}} & \textbf{\makecell{PLL\_F64M\_CLK\\64 MHz}} & \textbf{\makecell{PLL\_F48M\_CLK\\48 MHz}} & \multirow{-2}{*}{\textbf{\makecell{RC\_FAST\_CLK\\8 MHz}}} & \multirow{-2}{*}{\textbf{\makecell{RC\_SLOW\_CLK\\130 kHz}}} & \multirow{-2}{*}{\textbf{\makecell{OSC\_SLOW\_CLK\\32 kHz}}} & \multirow{-2}{*}{\textbf{\makecell{XTAL32K\_CLK\\32 kHz}}} & \multirow{-2}{*}{\textbf{\makecell{PLL\_LP\_CLK\\8 MHz}}} & \textbf{96 MHz/64 MHz/32 MHz/8 MHz} & \multirow{-2}{*}{\textbf{CRYPTO\_CLK}} & \multirow{-2}{*}{\textbf{APB\_CLK}} &  \multirow{-2}{*}{\textbf{AHB\_CLK}} & \multirow{-2}{*}{\textbf{CPU\_CLK}} & \multirow{-2}{*}{\textbf{\makecell{LP\_DYN\_\\FAST\_CLK}}}  & \multirow{-2}{*}{\textbf{\makecell{LP\_DYN\_\\SLOW\_CLK}}} & \multirow{-2}{*}{\textbf{\makecell{XTAL\_D2\_CLK\\16 MHz}}} &  \multirow{-2}{*}{\textbf{LP\_FAST\_CLK}} & \multirow{-2}{*}{\textbf{\makecell{Clock\\from IO}}} \\ \hline
    \endhead

    \multicolumn{20}{r}{\textbf{Continued on the next page...}} \\
    \endfoot
    \endlastfoot
    \hline

    PLL\_F48M\_CLK &   & Y &   &  &   &   &   &   &  &   &   &   &   &   &   &   &   &  &  \\ \hline
    HP\_ROOT\_CLK & Y & Y & Y &  & Y &   &   &   &  &   &   &   &   &   &   &   &   &  &  \\ \hline
    CRYPTO\_CLK & Y & Y & Y &  & Y &   &   &   &  &   &   &   &   &   &   &   &   &  &  \\ \hline
    APB\_CLK &   &   &   &  &   &   &   &   &  & Y &   &   &   &   &   &   &   &  &  \\ \hline
    AHB\_CLK &   &   &   &  &   &   &   &   &  & Y &   &   &   &   &   &   &   &  &  \\ \hline
    CPU\_CLK &   &   &   &  &   &   &   &   &  & Y &   &   &   &   &   &   &   &  &  \\ \hline
\end{longtable}

\begin{longtable}[c]{|m{2.1cm}|c|c|c|c|c|c|c|c|c|C{1.4cm}|c|c|c|c|c|c|c|c|c| }

 \caption{HP Clocks Used by Each Peripheral}
    \label{tab:resclk-periph-clks-hp} \\
  \hline
    \rowcolor{lightgray}
    & \multicolumn{3}{c|}{\textbf{Source Clock}} & \textbf{Derived Clock} & \multicolumn{5}{c|}{\textbf{Source Clock}} & \multicolumn{9}{c|}{\textbf{Derived Clock}} & \textbf{Source Clock} \\ \cline{2-20}
    \rowcolor{lightgray}
    & & \multicolumn{3}{c|}{\textbf{PLL\_CLK}} &  &  &  &  &  & \textbf{HP\_ROOT\_CLK} &  &  &  &  &  &  & &  & \multirow{2}{*}{\textbf{}} \\ \cline{3-5}
    \rowcolor{lightgray}
    \multirow{-3}{*}{\textbf{Peripheral}} & \multirow{-2}{*}{\textbf{\makecell{XTAL\_CLK\\32 MHz}}} & \textbf{\makecell{PLL\_F96M\_CLK\\96 MHz}} & \textbf{\makecell{PLL\_F64M\_CLK\\64 MHz}} & \textbf{\makecell{PLL\_F48M\_CLK\\48 MHz}} & \multirow{-2}{*}{\textbf{\makecell{RC\_FAST\_CLK\\8 MHz}}} & \multirow{-2}{*}{\textbf{\makecell{RC\_SLOW\_CLK\\130 kHz}}} & \multirow{-2}{*}{\textbf{\makecell{OSC\_SLOW\_CLK\\32 kHz}}} & \multirow{-2}{*}{\textbf{\makecell{XTAL32K\_CLK\\32 kHz}}} & \multirow{-2}{*}{\textbf{\makecell{PLL\_LP\_CLK\\8 MHz}}} & \textbf{96 MHz/64 MHz/32 MHz/8 MHz} & \multirow{-2}{*}{\textbf{CRYPTO\_CLK}} & \multirow{-2}{*}{\textbf{APB\_CLK}} &  \multirow{-2}{*}{\textbf{AHB\_CLK}} & \multirow{-2}{*}{\textbf{CPU\_CLK}} & \multirow{-2}{*}{\textbf{\makecell{LP\_DYN\_\\FAST\_CLK}}}  & \multirow{-2}{*}{\textbf{\makecell{LP\_DYN\_\\SLOW\_CLK}}} & \multirow{-2}{*}{\textbf{\makecell{XTAL\_D2\_CLK\\16 MHz}}} &  \multirow{-2}{*}{\textbf{LP\_FAST\_CLK}} & \multirow{-2}{*}{\textbf{\makecell{Clock\\from IO}}} \\ \hline
        \endfirsthead

    \multicolumn{20}{c}{\bfseries \tablename\ \thetable{} -- Continued from the previous page...}\\\hline
    \rowcolor{lightgray}
    & \multicolumn{3}{c|}{\textbf{Source Clock}} & \textbf{Derived Clock} & \multicolumn{5}{c|}{\textbf{Source Clock}} & \multicolumn{9}{c|}{\textbf{Derived Clock}} & \textbf{Source Clock} \\ \cline{2-20}
    \rowcolor{lightgray}
    & & \multicolumn{3}{c|}{\textbf{PLL\_CLK}} &  &  &  &  &  & \textbf{HP\_ROOT\_CLK} &  &  &  &  &  &  & &  & \multirow{2}{*}{\textbf{}} \\ \cline{3-5}
    \rowcolor{lightgray}
    \multirow{-3}{*}{\textbf{Peripheral}} & \multirow{-2}{*}{\textbf{\makecell{XTAL\_CLK\\32 MHz}}} & \textbf{\makecell{PLL\_F96M\_CLK\\96 MHz}} & \textbf{\makecell{PLL\_F64M\_CLK\\64 MHz}} & \textbf{\makecell{PLL\_F48M\_CLK\\48 MHz}} & \multirow{-2}{*}{\textbf{\makecell{RC\_FAST\_CLK\\8 MHz}}} & \multirow{-2}{*}{\textbf{\makecell{RC\_SLOW\_CLK\\130 kHz}}} & \multirow{-2}{*}{\textbf{\makecell{OSC\_SLOW\_CLK\\32 kHz}}} & \multirow{-2}{*}{\textbf{\makecell{XTAL32K\_CLK\\32 kHz}}} & \multirow{-2}{*}{\textbf{\makecell{PLL\_LP\_CLK\\8 MHz}}} & \textbf{96 MHz/64 MHz/32 MHz/8 MHz} & \multirow{-2}{*}{\textbf{CRYPTO\_CLK}} & \multirow{-2}{*}{\textbf{APB\_CLK}} &  \multirow{-2}{*}{\textbf{AHB\_CLK}} & \multirow{-2}{*}{\textbf{CPU\_CLK}} & \multirow{-2}{*}{\textbf{\makecell{LP\_DYN\_\\FAST\_CLK}}}  & \multirow{-2}{*}{\textbf{\makecell{LP\_DYN\_\\SLOW\_CLK}}} & \multirow{-2}{*}{\textbf{\makecell{XTAL\_D2\_CLK\\16 MHz}}} &  \multirow{-2}{*}{\textbf{LP\_FAST\_CLK}} & \multirow{-2}{*}{\textbf{\makecell{Clock\\from IO}}} \\ \hline
    \endhead

    \multicolumn{20}{r}{\textbf{Continued on the next page...}} \\
    \endfoot
    \endlastfoot
    \hline

    \nameref{mod:timg} & Y &  &  & Y & Y &  &  &  &  &  &  &  &  &  &  &  &  &  &  \\ \hline
    \hyperref[mod:wdt]{Main System Watchdog Timers (MWDT)} & Y &  &  & Y & Y &  &  &  &  &  &  &  &  &  &  &  &  &  &  \\ \hline
    \nameref{mod:i2s} & Y & Y & Y &  &  &  &  &  &  &  &  &  &  &  &  &  &  &  & I2S\_MCLK\_in \\ \hline
    \nameref{mod:uart} & Y &  &  & Y & Y &  &  &  &  &  &  &  &  &  &  &  &  &  &  \\ \hline
    \nameref{mod:rmt} & Y &  &  &  & Y &  &  &  &  &  &  &  &  &  &  &  &  &  &  \\ \hline
    \nameref{mod:mcpwm} & Y & Y &  &  & Y &  &  &  &  &  &  &  &  &  &  &  &  &  &  \\ \hline
    \nameref{mod:i2c} & Y &  &  &  & Y &  &  &  &  &  &  &  &  &  &  &  &  &  &  \\ \hline
    \hyperref[mod:spi]{SPI2} & Y &  &  & Y & Y &  &  &  &  &  &  &  &  &  &  &  &  &  &  \\ \hline
    \hyperref[mod:adc]{SAR ADC}  & Y & Y &  &  & Y &  &  &  &  &  &  &  &  &  &  &  &  &  &  \\ \hline
    \nameref{mod:usbserialjtag} &  &  &  & Y &  &  &  &  &  &  &  &  &  &  &  &  &  &  &  \\ \hline
    \nameref{mod:twai} & Y &  &  &  & Y &  &  &  &  &  &  &  &  &  &  &  &  &  &  \\ \hline
    \nameref{mod:ledpwm} & Y & Y &  &  & Y &  &  &  &  &  &  &  &  &  &  &  &  &  &  \\ \hline
    \nameref{mod:systimer} & Y &  &  &  & Y &  &  &  &  &  &  &  &  &  &  &  &  &  &  \\ \hline
    \nameref{mod:parlio} & Y & Y &  &  & Y &  &  &  &  &  &  &  &  &  &  &  &  &  & \makecell[l]{parl\_rx\_clk\_in\\parl\_tx\_clk\_in} \\ \hline
    \hyperref[mod:iomuxgpio]{IO MUX} & Y &  &  & Y & Y &  &  &  &  &  &  &  &  &  &  &  &  &  &  \\ \hline
    \makecell[l]{\nameref{mod:aes}\\\nameref{mod:sha}\\\nameref{mod:rsa}\\\nameref{mod:ecc}\\\nameref{mod:digsig}\\\nameref{mod:hmac}\\\nameref{mod:ecdsa}} &  &  &  &  &  &  &  &  &  &  & Y &  &  &  &  &  &  &  &  \\ \hline
    \nameref{mod:intmtrx} &  &  &  &  &  &  &  &  &  &  &  &  & Y &  &  &  &  &  &  \\ \hline
    \nameref{mod:pcnt} &  &  &  &  &  &  &  &  &  &  &  & Y &  &  &  &  &  &  &  \\ \hline
    \nameref{mod:evntaskmatrix} &  &  &  &  &  &  &  &  &  &  &  &  & Y &  &  &  &  &  &  \\ \hline
    \nameref{mod:dma} &  &  &  &  &  &  &  &  &  &  &  &  & Y &  &  &  &  &  &  \\ \hline
    \hyperref[mod:uart]{UHCI} &  &  &  &  &  &  &  &  &  &  &  &  & Y &  &  &  &  &  &  \\ \hline
    \nameref{mod:dbgassist} &  &  &  &  &  &  &  &  &  &  &  &  &  & Y &  &  &  &  &  \\ \hline
    \nameref{mod:riscvtracenc}  &  &  &  &  &  &  &  &  &  &  &  &  & Y & Y &  &  &  &  &  \\ \hline
    \hyperref[mod:riscvcpu]{Interrupt priority registers (INTPRI)} &  &  &  &  &  &  &  &  &  &  &  &  &  & Y &  &  &  &  &  \\ \hline
\end{longtable}

\begin{longtable}[c]{|m{1.6cm}|c|c|c|c|c|c|c|c|c|C{1.5cm}|c|c|c|c|c|c|c|c|c| }

 \caption{Derived LP Clock Source}
    \label{tab:resclk-periph-clks-root-lp} \\
  \hline
    \rowcolor{lightgray}
    & \multicolumn{3}{c|}{\textbf{Source Clock}} & \textbf{Derived Clock} & \multicolumn{5}{c|}{\textbf{Source Clock}} & \multicolumn{9}{c|}{\textbf{Derived Clock}} & \textbf{Source Clock} \\ \cline{2-20}
    \rowcolor{lightgray}
    & & \multicolumn{3}{c|}{\textbf{PLL\_CLK}} &  &  &  &  &  & \textbf{HP\_ROOT\_CLK} &  &  &  &  &  &  & &  & \multirow{2}{*}{\textbf{}} \\ \cline{3-5}
    \rowcolor{lightgray}
    \multirow{-3}{*}{\textbf{Derived Clock}} & \multirow{-2}{*}{\textbf{\makecell{XTAL\_CLK\\32 MHz}}} & \textbf{\makecell{PLL\_F96M\_CLK\\96 MHz}} & \textbf{\makecell{PLL\_F64M\_CLK\\64 MHz}} & \textbf{\makecell{PLL\_F48M\_CLK\\48 MHz}} & \multirow{-2}{*}{\textbf{\makecell{RC\_FAST\_CLK\\8 MHz}}} & \multirow{-2}{*}{\textbf{\makecell{RC\_SLOW\_CLK\\130 kHz}}} & \multirow{-2}{*}{\textbf{\makecell{OSC\_SLOW\_CLK\\32 kHz}}} & \multirow{-2}{*}{\textbf{\makecell{XTAL32K\_CLK\\32 kHz}}} & \multirow{-2}{*}{\textbf{\makecell{PLL\_LP\_CLK\\8 MHz}}} & \textbf{96 MHz/64 MHz/32 MHz/8 MHz} & \multirow{-2}{*}{\textbf{CRYPTO\_CLK}} & \multirow{-2}{*}{\textbf{APB\_CLK}} &  \multirow{-2}{*}{\textbf{AHB\_CLK}} & \multirow{-2}{*}{\textbf{CPU\_CLK}} & \multirow{-2}{*}{\textbf{\makecell{LP\_DYN\_\\FAST\_CLK}}}  & \multirow{-2}{*}{\textbf{\makecell{LP\_DYN\_\\SLOW\_CLK}}} & \multirow{-2}{*}{\textbf{\makecell{XTAL\_D2\_CLK\\16 MHz}}} &  \multirow{-2}{*}{\textbf{LP\_FAST\_CLK}} & \multirow{-2}{*}{\textbf{\makecell{Clock\\from IO}}} \\ \hline

        \endfirsthead

    \multicolumn{20}{c}{\bfseries \tablename\ \thetable{} -- Continued from the previous page...}\\\hline
    \rowcolor{lightgray}
    & \multicolumn{3}{c|}{\textbf{Source Clock}} & \textbf{Derived Clock} & \multicolumn{5}{c|}{\textbf{Source Clock}} & \multicolumn{9}{c|}{\textbf{Derived Clock}} & \textbf{Source Clock} \\ \cline{2-20}
    \rowcolor{lightgray}
    & & \multicolumn{3}{c|}{\textbf{PLL\_CLK}} &  &  &  &  &  & \textbf{HP\_ROOT\_CLK} &  &  &  &  &  &  & &  & \multirow{2}{*}{\textbf{}} \\ \cline{3-5}
    \rowcolor{lightgray}
    \multirow{-3}{*}{\textbf{Derived Clock}} & \multirow{-2}{*}{\textbf{\makecell{XTAL\_CLK\\32 MHz}}} & \textbf{\makecell{PLL\_F96M\_CLK\\96 MHz}} & \textbf{\makecell{PLL\_F64M\_CLK\\64 MHz}} & \textbf{\makecell{PLL\_F48M\_CLK\\48 MHz}} & \multirow{-2}{*}{\textbf{\makecell{RC\_FAST\_CLK\\8 MHz}}} & \multirow{-2}{*}{\textbf{\makecell{RC\_SLOW\_CLK\\130 kHz}}} & \multirow{-2}{*}{\textbf{\makecell{OSC\_SLOW\_CLK\\32 kHz}}} & \multirow{-2}{*}{\textbf{\makecell{XTAL32K\_CLK\\32 kHz}}} & \multirow{-2}{*}{\textbf{\makecell{PLL\_LP\_CLK\\8 MHz}}} & \textbf{96 MHz/64 MHz/32 MHz/8 MHz} & \multirow{-2}{*}{\textbf{CRYPTO\_CLK}} & \multirow{-2}{*}{\textbf{APB\_CLK}} &  \multirow{-2}{*}{\textbf{AHB\_CLK}} & \multirow{-2}{*}{\textbf{CPU\_CLK}} & \multirow{-2}{*}{\textbf{\makecell{LP\_DYN\_\\FAST\_CLK}}}  & \multirow{-2}{*}{\textbf{\makecell{LP\_DYN\_\\SLOW\_CLK}}} & \multirow{-2}{*}{\textbf{\makecell{XTAL\_D2\_CLK\\16 MHz}}} &  \multirow{-2}{*}{\textbf{LP\_FAST\_CLK}} & \multirow{-2}{*}{\textbf{\makecell{Clock\\from IO}}} \\ \hline
    \endhead

    \multicolumn{20}{r}{\textbf{Continued on the next page...}} \\
    \endfoot
    \endlastfoot
    \hline

    LP\_DYN\_FAST\_CLK &   &   &   &  &   & Y & Y & Y &  &   &   &   &   &   &   &   &   & Y &  \\ \hline
    LP\_DYN\_SLOW\_CLK &   &   &   &  &   & Y & Y & Y &  &   &   &   &   &   &   &   &   &   &  \\ \hline
    XTAL\_D2\_CLK & Y &   &   &  &   &   &   &   &  &   &   &   &   &   &   &   &   &  &  \\ \hline
    LP\_fAST\_CLK &   &   &   &  & Y &   &   &   & Y &   &   &   &   &   &   &   & Y &  &  \\ \hline
\end{longtable}

\begin{longtable}[c]{|m{2.1cm}|c|c|c|c|c|c|c|c|c|C{1.4cm}|c|c|c|c|c|c|c|c|c| }

 \caption{LP Clocks Used by Each Peripheral}
    \label{tab:resclk-periph-clks-lp} \\
  \hline
    \rowcolor{lightgray}
    & \multicolumn{3}{c|}{\textbf{Source Clock}} & \textbf{Derived Clock} & \multicolumn{5}{c|}{\textbf{Source Clock}} & \multicolumn{9}{c|}{\textbf{Derived Clock}} & \textbf{Source Clock} \\ \cline{2-20}
    \rowcolor{lightgray}
    & & \multicolumn{3}{c|}{\textbf{PLL\_CLK}} &  &  &  &  &  & \textbf{HP\_ROOT\_CLK} &  &  &  &  &  &  & &  & \multirow{2}{*}{\textbf{}} \\ \cline{3-5}
    \rowcolor{lightgray}
    \multirow{-3}{*}{\textbf{Peripheral}} & \multirow{-2}{*}{\textbf{\makecell{XTAL\_CLK\\32 MHz}}} & \textbf{\makecell{PLL\_F96M\_CLK\\96 MHz}} & \textbf{\makecell{PLL\_F64M\_CLK\\64 MHz}} & \textbf{\makecell{PLL\_F48M\_CLK\\48 MHz}} & \multirow{-2}{*}{\textbf{\makecell{RC\_FAST\_CLK\\8 MHz}}} & \multirow{-2}{*}{\textbf{\makecell{RC\_SLOW\_CLK\\130 kHz}}} & \multirow{-2}{*}{\textbf{\makecell{OSC\_SLOW\_CLK\\32 kHz}}} & \multirow{-2}{*}{\textbf{\makecell{XTAL32K\_CLK\\32 kHz}}} & \multirow{-2}{*}{\textbf{\makecell{PLL\_LP\_CLK\\8 MHz}}} & \textbf{96 MHz/64 MHz/32 MHz/8 MHz} & \multirow{-2}{*}{\textbf{CRYPTO\_CLK}} & \multirow{-2}{*}{\textbf{APB\_CLK}} &  \multirow{-2}{*}{\textbf{AHB\_CLK}} & \multirow{-2}{*}{\textbf{CPU\_CLK}} & \multirow{-2}{*}{\textbf{\makecell{LP\_DYN\_\\FAST\_CLK}}}  & \multirow{-2}{*}{\textbf{\makecell{LP\_DYN\_\\SLOW\_CLK}}} & \multirow{-2}{*}{\textbf{\makecell{XTAL\_D2\_CLK\\16 MHz}}} &  \multirow{-2}{*}{\textbf{LP\_FAST\_CLK}} & \multirow{-2}{*}{\textbf{\makecell{Clock\\from IO}}} \\ \hline
        \endfirsthead

    \multicolumn{20}{c}{\bfseries \tablename\ \thetable{} -- Continued from the previous page...}\\\hline
    \rowcolor{lightgray}
    & \multicolumn{3}{c|}{\textbf{Source Clock}} & \textbf{Derived Clock} & \multicolumn{5}{c|}{\textbf{Source Clock}} & \multicolumn{9}{c|}{\textbf{Derived Clock}} & \textbf{Source Clock} \\ \cline{2-20}
    \rowcolor{lightgray}
    & & \multicolumn{3}{c|}{\textbf{PLL\_CLK}} &  &  &  &  &  & \textbf{HP\_ROOT\_CLK} &  &  &  &  &  &  & &  & \multirow{2}{*}{\textbf{}} \\ \cline{3-5}
    \rowcolor{lightgray}
    \multirow{-3}{*}{\textbf{Peripheral}} & \multirow{-2}{*}{\textbf{\makecell{XTAL\_CLK\\32 MHz}}} & \textbf{\makecell{PLL\_F96M\_CLK\\96 MHz}} & \textbf{\makecell{PLL\_F64M\_CLK\\64 MHz}} & \textbf{\makecell{PLL\_F48M\_CLK\\48 MHz}} & \multirow{-2}{*}{\textbf{\makecell{RC\_FAST\_CLK\\8 MHz}}} & \multirow{-2}{*}{\textbf{\makecell{RC\_SLOW\_CLK\\130 kHz}}} & \multirow{-2}{*}{\textbf{\makecell{OSC\_SLOW\_CLK\\32 kHz}}} & \multirow{-2}{*}{\textbf{\makecell{XTAL32K\_CLK\\32 kHz}}} & \multirow{-2}{*}{\textbf{\makecell{PLL\_LP\_CLK\\8 MHz}}} & \textbf{96 MHz/64 MHz/32 MHz/8 MHz} & \multirow{-2}{*}{\textbf{CRYPTO\_CLK}} & \multirow{-2}{*}{\textbf{APB\_CLK}} &  \multirow{-2}{*}{\textbf{AHB\_CLK}} & \multirow{-2}{*}{\textbf{CPU\_CLK}} & \multirow{-2}{*}{\textbf{\makecell{LP\_DYN\_\\FAST\_CLK}}}  & \multirow{-2}{*}{\textbf{\makecell{LP\_DYN\_\\SLOW\_CLK}}} & \multirow{-2}{*}{\textbf{\makecell{XTAL\_D2\_CLK\\16 MHz}}} &  \multirow{-2}{*}{\textbf{LP\_FAST\_CLK}} & \multirow{-2}{*}{\textbf{\makecell{Clock\\from IO}}} \\ \hline
    \endhead

    \multicolumn{20}{r}{\textbf{Continued on the next page...}} \\
    \endfoot
    \endlastfoot
    \hline

    \hyperref[mod:efuse]{eFuse Controller (eFuse)} &  &  &  &  &  &  &  &  &  &  &  &  &  &  &  &  & Y &  &  \\ \hline
    \hyperref[mod:wdt]{RTC Watchdog Timer (RWDT)}  &  &  &  &  &  &  &  &  &  &  &  &  &  &  & Y & Y &  &  &  \\ \hline
    \hyperref[mod:lowpowm]{RTC Timer (RTC Timer)}  &  &  &  &  &  &  &  &  &  &  &  &  &  &  & Y & Y &  &  &  \\ \hline
    \hyperref[mod:lowpowm]{Brownout Detector} &  &  &  &  &  &  &  &  &  &  &  &  &  &  & Y &  &  &  &  \\ \hline
    \hyperref[mod:lowpowm]{Power Management Unit (PMU)} &  &  &  &  &  &  &  &  &  &  &  &  &  &  & Y & Y &  & Y &  \\ \hline
\end{longtable}
\end{tiny}

\end{landscape}

\textbf{PLL\_F96M\_CLK}

PLL\_F96M\_CLK is a 96 MHz clock. PLL\_F48M\_CLK (48 MHz) is divided from PLL\_F96M\_CLK.


\textbf{CRYPTO\_CLK}

As shown in Figure \ref{tab:resclk-periph-clks-root-hp}, CRYPTO\_CLK can be derived from XTAL\_CLK, PLL\_96M\_CLK, PLL\_64M\_CLK, or RC\_FAST\_CLK, and its frequency is up to 96 MHz.


To protect encryption and decryption peripherals from DPA (Differential Power Analysis) attacks, a random divider strategy is implemented for the functional clock of encryption and decryption peripherals. Three security levels are available, depending on the range of random divider. Users can select the security level by configuring  \hyperref[regdesc:HPSYSTEMSECDPACONFREG]{HP\_SYSTEM\_SEC\_DPA\_CONF\_REG}. If  \hyperref[fielddesc:HPSYSTEMSECDPACFGSEL]{HP\_SYSTEM\_SEC\_DPA\_CFG\_SEL} is set to 1, the security level is determined by the configuration of \hyperref[fielddesc:EFUSESECDPALEVEL]{EFUSE\_SEC\_DPA\_LEVEL}, otherwise, by the value of  \hyperref[fielddesc:HPSYSTEMSECDPALEVEL]{HP\_SYSTEM\_SEC\_DPA\_LEVEL}.


\textbf{LED\_PWM Clock}

LEDC module uses PLL\_F96M\_CLK, RC\_FAST\_CLK or XTAL\_CLK as its clock source. When the system is in low-power mode (APB\_CLK is disabled), most peripherals are halted, but LEDC can still work via RC\_FAST\_CLK.




\section{Programming Procedures}
\label{sec:resclk-prog-procedures}

\subsection{HP System Clock Configuration}

When configuring \hyperref[fielddesc:PCRSOCCLKSEL]{PCR\_SOC\_CLK\_SEL} to select the clock source of HP\_ROOT\_CLK, or configuring the clock divisor for CPU\_CLK via \hyperref[fielddesc:PCRCPUDIVNUM]{PCR\_CPU\_DIV\_NUM} and AHB\_CLK via \hyperref[fielddesc:PCRAHBDIVNUM]{PCR\_AHB\_DIV\_NUM}, please also set the enable register to apply the new configuration. To check whether the new configuration takes effect, read \hyperref[fielddesc:PCRBUSCLOCKUPDATE]{PCR\_BUS\_CLOCK\_UPDATE} and see if it is 0.

\subsection{LP System Clock Configuration}

The clock source of LP\_SLOW\_CLK can be configured via  \hyperref[fielddesc:LPCLKRSTSLOWCLKSEL]{LP\_CLKRST\_SLOW\_CLK\_SEL}.

The clock source of LP\_FAST\_CLK can be configured via  \hyperref[fielddesc:LPCLKRSTFASTCLKSEL]{LP\_CLKRST\_FAST\_CLK\_SEL}.

\subsection{Peripheral Clock Reset and Configuration}\label{sec:resclk-periperal-clock-configuration}

\begin{importantlisting}
\vspace{-0.5em}

\chipname{} features low power consumption. This is why some peripheral clocks are gated (disabled) by default. Before using any of these peripherals, it is mandatory to enable the clock for the given peripheral by setting the corresponding CLK\_EN bit to 1, and release the peripheral from reset state to make it operational by setting the RST\_EN bit to 0.
\end{importantlisting}

The clocks of most peripherals can be classified into two types:
\begin{itemize}
    \item Bus clock: used to configure peripheral registers.
    \item Functional clock: such as UART's reference clock, used by peripherals to operate.
\end{itemize}
The functional clock of most peripherals can be selected from multiple clock sources. For clock gating registers, whether they are used to gate the bus clock (AHB\_CLK, APB\_CLK) or the functional clock will be stated in the corresponding register descriptions.

Bus clock switches, functional clock switches, and configuration registers for clock source selection and clock frequency division are grouped into the PCR module. For more information, see Section \ref{sec:resclk-reg-summ} \textit{\nameref{sec:resclk-reg-summ}}.

When a peripheral is not working, users can turn off its functional clock by configuring related PCR registers. Turning off the peripheral's functional clock does not affect the rest of the system.

Take the I2C clock configuration as an example.


\begin{figure}[H]
    \centering
    \includegraphics[width=0.8\textwidth]{\modulefiles/figures/sys_clk_struct_pcr.eps}
    \caption{Clock Configuration Example}
    \label{figure:Clock Config}
\end{figure}


Figure \ref{figure:Clock Config} shows the clock structure of I2C. The clock structure of other peripherals is similar to this one. CLK\_SWITCH is used to select a clock output and CLK\_GATE to turn on/off the clock.

In scenarios that require low power consumption, when the peripheral is not in use, in addition to turning off the functional clock, the bus clock of the peripheral can also be turned off to further lower the power consumption.

Note that if you turn off the bus clock first, the functional clock may continue working. Therefore, when turning off clocks, it is recommended to turn off the functional clock first and then the bus clock; when turning on clocks, it is recommended to  turn on the bus clock first and then the functional clock.

\begin{tiplisting}
In this chapter, all divisor configuration registers are configured with the actual divisor minus 1.
\end{tiplisting}

% >>> If your module has no registers, delete sample text
%       for Sections Register Summary and Registers below <<<

%%%%%%%%%%%%%%%%%%%%%%%%
%%% Register Summary %%%
%%%%%%%%%%%%%%%%%%%%%%%%

% >>> Update the sample text below and insert
%     the info on register summary into the subfile <<<

\clearpage
\hypertarget{resclk-reg-summ}{}
\section{Register Summary}
\label{sec:resclk-reg-summ}

\subsection{PCR Register Summary}

The addresses in this section are relative to Power/Clock/Reset (PCR) Register base address provided in Table \ref{tab:sysmem-base-address} in Chapter \ref{mod:sysmem} \textit{\nameref{mod:sysmem}}.

The abbreviations given in Column \textbf{Access} are explained in Section \hyperref[glossary-access-types]{Access Types for Registers}.

\subfile{\modulefiles/08-PCR-reg-summ__EN}

\subsection{LP System Clock Register Summary}

The addresses of the last two registers with the LPPERI prefix in this section are relative to the Low-power Peripheral Register (LPPERI) base address. The addresses of the third and fourth to last registers with the LP\_AON prefix in this section are relative to the Low-power Always-on Register (LP\_AON) base address. The other addresses in this section are relative to the Low-power Clock/Reset Register (LP\_CLKRST) base address. For base address, please refer to Table \ref{tab:sysmem-base-address} in Chapter \ref{mod:sysmem} \textit{\nameref{mod:sysmem}}.

The abbreviations given in Column \textbf{Access} are explained in Section \hyperref[glossary-access-types]{Access Types for Registers}.

\subfile{\modulefiles/08-LP-CLKRST-reg-summ__EN}


%%%%%%%%%%%%%%%%%
%%% Registers %%%
%%%%%%%%%%%%%%%%%

% >>> Update the sample text below and insert
%      the info on registers into the subfile <<<

\clearpage
\section{Registers}

\subsection{PCR Registers}


The addresses in this section are relative to the Power/Clock/Reset (PCR) Register base address provided in Table  \ref{tab:sysmem-base-address} in Chapter  \ref{mod:sysmem} \textit{\nameref{mod:sysmem}}.


\subfile{\modulefiles/08-PCR-reg__EN}


\subsection{LP System Clock Registers}


The addresses of the last two registers with the LPPERI prefix in this section are relative to the Low-power Peripheral Register (LPPERI) base address.  The addresses of the third and fourth to last registers with the LP\_AON prefix in this section are relative to the Low-power Always-on Register (LP\_AON) base address.The other addresses in this section are relative to the Low-power Clock/Reset Register (LP\_CLKRST) base address. For base address, please refer to Table  \ref{tab:sysmem-base-address} in Chapter \ref{mod:sysmem} \textit{\nameref{mod:sysmem}}.


\subfile{\modulefiles/08-LP-CLKRST-reg__EN}

\end{document}
